\section{Related Work}

The application of Large Language Models (LLMs) to forecasting has rapidly evolved, branching into numerical and qualitative domains.

\para{Numerical Time Series Forecasting.} Early efforts in deep learning for time series focused on architectures like Informer and Autoformer. Recently, LLMs have emerged as powerful alternatives. \timellm \cite{jin2024timellm} introduces a reprogramming framework that maps time-series patches into the text space, allowing frozen LLMs to perform forecasting. Similarly, \onefitsall \cite{zhou2023onefitsall} leverages pretrained language models to achieve state-of-the-art results on several benchmarks. \logollm \cite{ou2025logollm} further enhances this by integrating local and global modeling. These works primarily focus on the high-data regime where patterns are buried in large numerical arrays. Our work differs by evaluating the zero-shot performance of foundational models without such specialized reprogramming, providing a baseline for generalist model capability.

\para{Event-based Forecasting.} Long-term forecasting of real-world events has traditionally been the domain of human intelligence. Superforecasters, individuals who consistently outperform the average in probabilistic estimation, have been the primary benchmark for this task. \forecastbench \cite{karger2025forecastbench} provides a dynamic framework for assessing AI models on such questions. Recent studies by \citet{lu2025evaluating} suggest that while frontier LLMs are surpassing the human public crowd, they still lag behind elite human superforecasters in many scenarios. Our findings contribute to this ongoing debate by showing that in certain event-based subsets, models like \methodname can indeed match or exceed expert performance.

\para{Pitfalls and Evaluation.} Evaluating LLM forecasters is fraught with challenges. \citet{paleka2025pitfalls} highlight the risk of temporal leakage, where a model's training data might already contain information about the events it is asked to predict. Additionally, LLMs often suffer from calibration issues, manifesting as overconfidence. Our study acknowledges these limitations and employs standard metrics like the Brier score to assess probabilistic accuracy while maintaining a critical view of potential data leakage in recent benchmarks.
